% =============================================================
% 5 Анализ радиального подшипника скольжения
%   центробежного насоса
% =============================================================
\chapter{Анализ радиального подшипника скольжения центробежного насоса}

\section{Объект исследования и исходные данные}

Объектом исследования в настоящей главе является радиальный
подшипник скольжения ротора центробежного насоса.
Подшипниковый узел предназначен для восприятия радиальных
нагрузок, передаваемых от ротора к корпусу, и обеспечения
жидкостного режима смазки в рабочем диапазоне скоростей
и нагрузок.

Подшипник представляет собой цилиндрическую пару: вал
радиуса~$R$ вращается внутри вкладыша с радиальным зазором~$c$.
Относительный зазор $\psi = c/R$ определяет геометрическую
характеристику подшипника; отношение длины к диаметру~$L/D$
характеризует его осевую протяжённость. Рабочей средой является
масло с динамической вязкостью~$\mu_0$ при референсной
температуре~$T_0$. Вал вращается с угловой скоростью~$\omega$
и нагружен радиальной силой~$F_{\mathrm{ext}}$, действующей
в плоскости поперечного сечения.

Расчёт выполняется в изотермическом приближении: температура
смазки принята постоянной ($T = T_0$, $\mu = \mathrm{const}$),
что обосновано в разделе~4.2 при умеренных скоростях скольжения
и нагрузках, характерных для рассматриваемого узла. Плотность
смазки считается постоянной ($\rho = \mathrm{const}$),
поверхности~-- гладкими (обоснование~-- раздел~4.3).

Основные параметры подшипника сведены в таблицу~\ref{tab:journal_params}.

\begin{table}[!ht]
\centering
\caption{Параметры радиального подшипника скольжения}
\label{tab:journal_params}
\begin{tabular}{lccc}
\hline
Параметр & Обозначение & Значение & Единица \\
\hline
Радиус вала & $R$ & --- & мм \\
Длина подшипника & $L$ & --- & мм \\
Радиальный зазор & $c$ & --- & мкм \\
Относительный зазор & $\psi = c/R$ & --- & -- \\
Отношение $L/D$ & $L/D$ & --- & -- \\
Угловая скорость & $\omega$ & --- & рад/с \\
Динамическая вязкость & $\mu_0$ & --- & Па$\cdot$с \\
Давление кавитации & $p_{\mathrm{cav}}$ & --- & Па \\
\hline
\end{tabular}
\end{table}

Схема подшипника в поперечном сечении представлена на
рисунке~\ref{fig:journal_scheme}.

\begin{figure}[!ht]
\centering
\fbox{\parbox{0.85\textwidth}{\centering
  \vspace{3.5cm}
  ПЛЕЙСХОЛДЕР ДЛЯ РИСУНКА\\
  Поперечное сечение радиального подшипника:\\
  вал (радиус $R$), вкладыш, зазор $h(\theta)$,\\
  эксцентриситет $e$, направление вращения $\omega$,\\
  внешняя нагрузка $F_{\mathrm{ext}}$, угол нагружения $\psi$
  \vspace{3.5cm}
}}
\caption{Схема радиального подшипника скольжения центробежного насоса}
\label{fig:journal_scheme}
\end{figure}


\section{Постановка задачи}

Расчётная область представляет собой развёрнутую поверхность
смазочного зазора $(\theta, z) \in [0,\,2\pi] \times [0,\,L]$,
где $\theta$~-- окружная координата, $z$~-- осевая координата.
При переходе к безразмерным переменным (подраздел~2.1.6)
осевая координата нормируется на полудлину: $Z = z/L$.

Для гладкого цилиндрического подшипника безразмерная толщина
зазора определяется законом~\eqref{eq:journal_h}:
\begin{equation}
  H(\theta) = 1 + \varepsilon\cos\theta,
  \label{eq:journal_h_ch5}
\end{equation}
где $\varepsilon = e/c$~-- относительный эксцентриситет
($0 \leq \varepsilon < 1$), $e$~-- смещение оси вала
относительно оси вкладыша. При $\varepsilon = 0$ зазор
равномерный, при $\varepsilon \to 1$ минимальный зазор
$h_{\min} = c\,(1 - \varepsilon)$ стремится к нулю.

Для подшипника с детерминированными микроструктурами толщина
зазора записывается в виде суперпозиции~\eqref{eq:h_superposition}:
\begin{equation}
  H = H_{\mathrm{base}}(\theta) + \Delta H_{\mathrm{texture}}(\theta, Z),
  \label{eq:journal_h_tex_ch5}
\end{equation}
где $H_{\mathrm{base}} = 1 + \varepsilon\cos\theta$~-- базовый зазор
гладкого подшипника, $\Delta H_{\mathrm{texture}}$~-- добавка,
описывающая микрорельеф (подраздел~2.2, формула~\eqref{eq:journal_h_textured}).

Распределение давления в смазочном слое описывается безразмерным
уравнением Рейнольдса~\eqref{eq:reynolds_nondim} с параметром
анизотропии~$\Lambda$~\eqref{eq:nondim_Lambda}. Для радиального
подшипника $\Lambda = 2R/L$; данный параметр определяет
соотношение окружного и осевого масштабов задачи.

Кавитация учитывается в рамках постановки JFO
(\eqref{eq:jfo_conservation}, \eqref{eq:jfo_complementarity}):
в зоне кавитации давление фиксируется на уровне
$P = 0$ (что соответствует $p = p_{\mathrm{cav}}$ в размерных
переменных), а вместо уравнения Рейнольдса решается уравнение
сохранения массы для степени заполнения~$\vartheta$.

Граничные условия задачи формулируются следующим образом:
по окружной координате~$\theta$ поле давления периодично
($P(\theta = 0) = P(\theta = 2\pi)$); на торцах подшипника
давление равно атмосферному ($P = 0$ при $Z = 0$ и $Z = 1$).

Допущения, принятые в настоящей главе, включают: изотермическое
приближение ($T = T_0$, обоснование~-- раздел~4.2), постоянство
плотности ($\rho = \mathrm{const}$, раздел~4.1) и гладкость
поверхностей ($\Delta h_{\mathrm{rough}} = 0$, раздел~4.3).

Расчётная область с граничными условиями схематически показана
на рисунке~\ref{fig:journal_domain}.

\begin{figure}[!ht]
\centering
\fbox{\parbox{0.85\textwidth}{\centering
  \vspace{3.0cm}
  ПЛЕЙСХОЛДЕР ДЛЯ РИСУНКА\\
  Развёрнутая расчётная область $(\theta, Z)$:\\
  периодичность по $\theta$, $P = 0$ на торцах $Z=0$ и $Z=1$,\\
  зона кавитации (схематически)
  \vspace{3.0cm}
}}
\caption{Расчётная область и граничные условия}
\label{fig:journal_domain}
\end{figure}


\section{Результаты для гладкого подшипника}

Результаты расчёта гладкого подшипника (без микроструктур)
составляют эталонную базу для последующего сравнения
с текстурированным вариантом.

\textit{Поле давления и кавитация.}
При стационарном вращении вала под нагрузкой ось вала смещается
на величину эксцентриситета~$e$ относительно оси вкладыша.
В сходящейся части зазора ($0 < \theta < \pi$ при выбранной
системе координат) смазка увлекается в область с уменьшающимся
зазором, что формирует зону повышенного давления. В расходящейся
части ($\pi < \theta < 2\pi$) зазор увеличивается; при
отсутствии кавитации давление здесь становилось бы отрицательным
(ниже давления насыщения). Постановка JFO
(\eqref{eq:jfo_conservation}, \eqref{eq:jfo_complementarity})
обнуляет давление в области, где оно падает ниже давления
кавитации, и формирует самосогласованную границу кавитационной
зоны, в которой смазочный слой заполнен не полностью
($\vartheta < 1$). Распределение давления по окружной
и осевой координатам, а также расположение кавитационной зоны
показаны на рисунке~\ref{fig:smooth_pressure}.

\begin{figure}[!ht]
\centering
\fbox{\parbox{0.85\textwidth}{\centering
  \vspace{3.5cm}
  ПЛЕЙСХОЛДЕР ДЛЯ РИСУНКА\\
  Слева: поле безразмерного давления $P(\theta, Z)$ для гладкого подшипника\\
  (2D-карта или 3D-поверхность, $\varepsilon = \varepsilon_{\mathrm{nom}}$).\\
  Справа: карта кавитации~-- зона $\vartheta < 1$ (где $P = 0$)
  \vspace{3.5cm}
}}
\caption{Поле давления и кавитационная зона гладкого подшипника}
\label{fig:smooth_pressure}
\end{figure}

\textit{Интегральные характеристики.}
По полученному полю давления вычисляются интегральные
характеристики подшипника: компоненты несущей силы
и результирующая сила~$F$ с углом нагружения~$\psi$
(\eqref{eq:func_F_components}, \eqref{eq:func_psi}),
безразмерная сила трения~$f$ и коэффициент трения
$\mu_f = f/F$, а также осевые утечки~$Q_{\mathrm{out}}$
(\eqref{eq:func_Q_journal}). Значения указанных характеристик
при номинальном эксцентриситете приведены в
таблице~\ref{tab:smooth_results}.

\begin{table}[!ht]
\centering
\caption{Интегральные характеристики гладкого подшипника
  при $\varepsilon = \varepsilon_{\mathrm{nom}}$}
\label{tab:smooth_results}
\begin{tabular}{lcc}
\hline
Характеристика & Обозначение & Значение \\
\hline
Несущая сила (безразмерная) & $F$ & --- \\
Угол нагружения & $\psi$ & --- \\
Коэффициент трения & $\mu_f$ & --- \\
Осевые утечки (безразмерные) & $Q_{\mathrm{out}}$ & --- \\
\hline
\end{tabular}
\end{table}

\textit{Зависимости от эксцентриситета.}
Характерное поведение интегральных характеристик
гидродинамического подшипника при изменении эксцентриситета
определяется нелинейной зависимостью толщины зазора от~$\varepsilon$.
Несущая способность~$F(\varepsilon)$ монотонно возрастает
с ростом эксцентриситета, причём при $\varepsilon \to 1$
рост приобретает резко нелинейный характер, обусловленный
сужением минимального зазора и соответствующим увеличением
пиковых давлений. Коэффициент трения~$\mu_f(\varepsilon)$
убывает с ростом~$\varepsilon$: хотя абсолютные потери на трение
растут из-за повышения вязкого сдвига в зоне минимального зазора,
несущая способность растёт быстрее, и отношение~$\mu_f = f/F$
уменьшается. Осевые утечки~$Q_{\mathrm{out}}(\varepsilon)$
убывают с ростом~$\varepsilon$, поскольку с увеличением
эксцентриситета расширяется кавитационная зона, ограничивающая
область сквозного осевого потока. Указанные тенденции
являются типичными для гидродинамических подшипников
и подтверждаются классическими результатами. Графики
зависимостей~$F(\varepsilon)$, $\mu_f(\varepsilon)$
и~$Q_{\mathrm{out}}(\varepsilon)$ представлены
на рисунке~\ref{fig:smooth_curves}.

\begin{figure}[!ht]
\centering
\fbox{\parbox{0.85\textwidth}{\centering
  \vspace{3.5cm}
  ПЛЕЙСХОЛДЕР ДЛЯ РИСУНКА\\
  Три графика зависимости от $\varepsilon$:\\
  $F(\varepsilon)$ (несущая сила),\\
  $\mu_f(\varepsilon)$ (коэффициент трения),\\
  $Q_{\mathrm{out}}(\varepsilon)$ (утечки).\\
  Кривые для гладкого подшипника.
  \vspace{3.5cm}
}}
\caption{Зависимость интегральных характеристик гладкого подшипника
  от эксцентриситета}
\label{fig:smooth_curves}
\end{figure}

Полученные результаты для гладкого подшипника служат эталоном,
относительно которого оценивается влияние детерминированных
микроструктур в последующих разделах.


\section{Выбор и параметры микроструктуры}

В качестве элементов микрорельефа выбраны эллипсоидальные лунки
с профилем~\eqref{eq:profile_ellipsoidal}, представляющие
наиболее распространённый тип текстуры в задачах
гидродинамической смазки. Раскладка элементов выполнена
по шахматной схеме~\eqref{eq:layout_staggered}, обеспечивающей
равномерное заполнение текстурированной зоны при заданном
коэффициенте покрытия.

Выбор зоны нанесения микроструктур определяется физическим
механизмом их воздействия на поле давления. В радиальном
подшипнике текстурирование может охватывать всю поверхность
вкладыша, только зону сходящегося клина, зону максимального
давления или зону расходящегося клина. Расположение текстурированной
области относительно зоны давления существенно влияет на результат:
лунки в области сходящегося зазора способны усилить
гидродинамический эффект, тогда как лунки в зоне расходящегося
зазора могут расширить кавитационную область.

Безразмерные параметры текстуры определены в подразделе~2.2.4:
глубина элемента~$H_p$~\eqref{eq:Hp_nondim}, безразмерные
полуоси~$A^*$, $B^*$ и коэффициент покрытия~$\varphi$.
Рассматриваются три конфигурации текстуры, различающиеся
глубиной элементов, плотностью покрытия и зоной нанесения
(таблица~\ref{tab:texture_configs}).

\begin{table}[!ht]
\centering
\caption{Параметры рассматриваемых конфигураций текстуры}
\label{tab:texture_configs}
\begin{tabular}{lccccc}
\hline
Конфигурация & $H_p$ & $A^*$ & $B^*$ & $\varphi$ & Зона \\
\hline
T1 & --- & --- & --- & --- & --- \\
T2 & --- & --- & --- & --- & --- \\
T3 & --- & --- & --- & --- & --- \\
\hline
\end{tabular}
\end{table}

Схема текстурированной поверхности на развёрнутой области
представлена на рисунке~\ref{fig:texture_layout}.

\begin{figure}[!ht]
\centering
\fbox{\parbox{0.85\textwidth}{\centering
  \vspace{3.0cm}
  ПЛЕЙСХОЛДЕР ДЛЯ РИСУНКА\\
  Развёрнутая поверхность $(\theta, Z)$ с зоной текстурирования:\\
  эллипсоидальные лунки в шахматной раскладке,\\
  обозначить $A^*$, $B^*$, шаги раскладки, зону нанесения
  \vspace{3.0cm}
}}
\caption{Схема зоны текстурирования на поверхности
  радиального подшипника}
\label{fig:texture_layout}
\end{figure}


\section{Сравнение гладкого и текстурированного подшипников}

\textit{Физический механизм влияния микроструктур.}
Эллипсоидальные лунки, нанесённые на рабочую поверхность
вкладыша, создают локальные расширения и сужения зазора,
формируя дополнительные микроклинья. В зоне сходящегося
зазора подшипника, где давление формируется за счёт
макроскопического клинового эффекта~\eqref{eq:journal_h_ch5},
лунки создают дополнительные локальные области конвергенции,
что может усилить генерацию давления и повысить несущую
способность. В расходящейся зоне, где давление снижается
до уровня кавитации, лунки способны расширить кавитационную
область, изменяя баланс потоков на границе фазового перехода.
Результирующее влияние микроструктур на интегральные
характеристики определяется совокупностью этих эффектов
и зависит от зоны нанесения, глубины и плотности покрытия
элементов.

\textit{Нормированные показатели эффективности.}
Для количественного сравнения текстурированного
и гладкого подшипников вводятся нормированные показатели:
\begin{equation}
  G_F = \frac{F_{\mathrm{tex}}}{F_{\mathrm{smooth}}}, \quad
  G_f = \frac{\mu_{f,\mathrm{tex}}}{\mu_{f,\mathrm{smooth}}}, \quad
  G_Q = \frac{Q_{\mathrm{out,tex}}}{Q_{\mathrm{out,smooth}}}, \quad
  G_h = \frac{h_{\min,\mathrm{tex}}}{h_{\min,\mathrm{smooth}}}.
  \label{eq:gain_factors}
\end{equation}
Значение~$G > 1$ указывает на рост соответствующего показателя
при текстурировании, $G < 1$~-- на снижение. С точки зрения
эксплуатации желательны~$G_F > 1$ (рост несущей способности),
$G_f < 1$ (снижение трения) и $G_h > 1$ (увеличение минимального
зазора, повышение запаса по толщине плёнки).

\textit{Поля давления и кавитация.}
Текстурирование поверхности вкладыша модифицирует распределение
давления в смазочном слое. В области расположения лунок поле
давления приобретает локальные возмущения~-- каждый элемент
микрорельефа формирует собственный микроклин, вносящий вклад
в интегральный баланс. Одновременно изменяется конфигурация
кавитационной зоны: форма границы кавитации и площадь области
с $\vartheta < 1$ зависят от параметров текстуры. Сопоставление
полей давления и карт кавитации для гладкого
и текстурированного подшипников при одинаковом значении
эксцентриситета представлено на рисунке~\ref{fig:comparison_pressure}.

\begin{figure}[!ht]
\centering
\fbox{\parbox{0.85\textwidth}{\centering
  \vspace{3.5cm}
  ПЛЕЙСХОЛДЕР ДЛЯ РИСУНКА\\
  Слева: $P(\theta, Z)$ для гладкого (верх) и текстурированного (низ) подшипников.\\
  Справа: карты кавитации для обоих случаев.\\
  Одинаковый $\varepsilon = \varepsilon_{\mathrm{nom}}$.
  \vspace{3.5cm}
}}
\caption{Поля давления и кавитационные зоны: гладкий
  и текстурированный подшипники}
\label{fig:comparison_pressure}
\end{figure}

\textit{Нагрузочные характеристики.}
Зависимости несущей способности~$F(\varepsilon)$
и коэффициента трения~$\mu_f(\varepsilon)$ для гладкого
и текстурированного подшипников позволяют оценить эффективность
текстурирования во всём диапазоне рабочих эксцентриситетов.
При малых эксцентриситетах ($\varepsilon \ll 1$), когда зазор
близок к равномерному, влияние микроструктур наиболее выражено:
каждая лунка вносит вклад, сопоставимый с макроскопическим
клиновым эффектом. С ростом~$\varepsilon$ собственный
гидродинамический эффект клинового зазора подшипника усиливается
и начинает доминировать, а относительный вклад текстуры
снижается. Графики нагрузочных характеристик для обоих вариантов
представлены на рисунке~\ref{fig:comparison_curves}.

\begin{figure}[!ht]
\centering
\fbox{\parbox{0.85\textwidth}{\centering
  \vspace{3.5cm}
  ПЛЕЙСХОЛДЕР ДЛЯ РИСУНКА\\
  Два графика (одна ось $\varepsilon$):\\
  $F(\varepsilon)$~-- гладкий и текстурированный (конфигурация T1),\\
  $\mu_f(\varepsilon)$~-- гладкий и текстурированный.\\
  Показать прирост несущей способности и снижение трения.
  \vspace{3.5cm}
}}
\caption{Сравнение нагрузочных характеристик гладкого
  и текстурированного подшипников}
\label{fig:comparison_curves}
\end{figure}


\section{Выводы по главе~5}

Проведён статический анализ радиального подшипника скольжения
центробежного насоса в изотермическом приближении с учётом
кавитации в рамках постановки JFO.

Для гладкого подшипника получены поля давления и установлены
границы кавитационной зоны при различных значениях
эксцентриситета. Зависимости несущей
способности~$F(\varepsilon)$, коэффициента
трения~$\mu_f(\varepsilon)$ и осевых
утечек~$Q_{\mathrm{out}}(\varepsilon)$ демонстрируют
характерное поведение гидродинамического подшипника: рост
несущей способности, убывание коэффициента трения и снижение
утечек с увеличением~$\varepsilon$.

Текстурирование поверхности вкладыша эллипсоидальными лунками
в шахматной раскладке модифицирует поле давления: в зоне
расположения элементов возникают локальные возмущения давления,
формируемые микроклиновым эффектом каждой лунки. Конфигурация
кавитационной зоны при текстурировании отличается от эталонного
случая гладкого подшипника.

Сравнение интегральных характеристик по нормированным
показателям~$G_F$, $G_f$, $G_Q$, $G_h$~\eqref{eq:gain_factors}
показывает, что при выбранных параметрах текстуры наблюдается
рост несущей способности ($G_F > 1$) и снижение коэффициента
трения ($G_f < 1$), наиболее выраженные при малых и умеренных
эксцентриситетах.

Эффективность текстурирования существенно зависит от зоны
нанесения микроструктур: расположение лунок в области
сходящегося клина оказывает более благоприятное влияние
на несущую способность, чем текстурирование расходящейся зоны
или полное покрытие поверхности. Глубина элементов~$H_p$
и коэффициент покрытия~$\varphi$ должны быть согласованы
с рабочим диапазоном эксцентриситетов.

При больших эксцентриситетах ($\varepsilon \to 1$) влияние
микроструктур ослабевает, поскольку макроскопический клиновой
эффект доминирует; в этом режиме текстурирование оказывается
нейтральным или малоэффективным.

Анализ радиального подшипника для иного класса объектов~--
опоры шарошечного долота~-- выполнен в главе~6. Упорный
подшипник рассмотрен в главе~7. Влияние микроструктур
на динамические характеристики подшипника (коэффициенты
жёсткости и демпфирования, устойчивость ротора) составляет
предмет главы~8.
